% *** ก่อนส่ง: ลบข้อความตัวอย่าง (\ldots) ทั้งหมด แล้วเขียนเนื้อหาของตนเอง ***
% =============================================================
% บทที่ 2 - ทฤษฎีและงานวิจัยที่เกี่ยวข้อง
% แก้ไขเนื้อหาในบทนี้ได้ที่ไฟล์ chapter2.tex
% =============================================================

% ทฤษฎีที่นำมาใช้ในการดำเนินโครงงาน
% อภิปรายงานวิจัย โครงงาน หรือระบบงานที่มีลักษณะเกี่ยวข้องหรือคล้ายคลึงกับโครงงานที่ทำ
% แสดงส่วนที่คล้ายกันและส่วนที่แตกต่างจากโครงงานนี้ เช่น จุดเด่นและข้อจำกัด แนวทางในการนำมาพัฒนาต่อ
% ใส่อ้างอิงให้ครบถ้วน

% -----------------------------------------------------------
\section{ทฤษฎีที่เกี่ยวข้อง}
% -----------------------------------------------------------

\subsection{ชื่อทฤษฎีหรือหลักการที่ 1}

อธิบายทฤษฎีหรือหลักการที่เกี่ยวข้อง\ldots{}~\cite{fielding2000}

\subsection{ชื่อทฤษฎีหรือหลักการที่ 2}

อธิบายทฤษฎีหรือหลักการที่เกี่ยวข้อง\ldots{}~\cite{cocalc2024}

% -----------------------------------------------------------
\section{งานวิจัยที่เกี่ยวข้อง}
% -----------------------------------------------------------

ในส่วนนี้จะนำเสนองานวิจัย โครงงาน หรือระบบงานที่มีลักษณะคล้ายคลึงกับโครงงานนี้ โดยจะกล่าวถึงแนวคิด วิธีการ รวมถึงจุดเด่นและข้อจำกัดของแต่ละงาน เพื่อนำมาเป็นแนวทางในการพัฒนา

\subsection{ชื่องานวิจัยที่ 1}

% อ้างอิงงานวิจัยด้วย \cite{key}
ผู้วิจัย \cite{example2024} ได้นำเสนอระบบ\ldots{} โดยใช้วิธีการ\ldots{} ผลการทดลองพบว่า\ldots{} อย่างไรก็ตาม งานวิจัยนี้มีข้อจำกัดในด้าน\ldots{}

\subsection{ชื่องานวิจัยที่ 2}

ผู้วิจัย \cite{example2023conf} ได้พัฒนา\ldots{} โดยใช้่วิธีการ\ldots{} ซึ่งมีแนวคิดคล้ายกับโครงงานนี้ในส่วนที่\ldots{} แต่ยังขาดการแก้ปัญหา\ldots{} ซึ่งโครงงานนี้นำเสนอการแก้่ปัญหาดังกล่าวโดย\ldots{}

\subsection{สรุปผลการเปรียบเทียบงานวิจัยที่เกี่ยวข้อง}

จากการทบทวนวรรณกรรมข้างต้น สามารถสรุปผลการเปรียบเทียบได้ดังตารางที่ \ref{tab:compare_research}

\begin{table}[H]
\centering
\caption{เปรียบเทียบงานวิจัยที่เกี่ยวข้อง}
\label{tab:compare_research}
\begin{tabular}{|l|c|c|c|c|}
\hline
\textbf{งานวิจัย} & \textbf{คุณสมบัติ A} & \textbf{คุณสมบัติ B} & \textbf{คุณสมบัติ C} & \textbf{คุณสมบัติ D} \\
\hline
งานวิจัยที่ 1 \cite{example2024} & \cmark & \cmark & -- & -- \\ \hline
งานวิจัยที่ 2 \cite{example2023conf} & \cmark & -- & \cmark & -- \\ \hline
\textbf{โครงงานนี้} & \cmark & \cmark & \cmark & \cmark \\
\hline
\end{tabular}
\end{table}

% ตัวอย่างการ cite หลายรูปแบบ:
% \cite{key1}            → [1]
% \cite{key1, key2}      → [1, 2]
% \cite{key1, key2, key3} → [1, 2, 3]

จากตารางที่ \ref{tab:compare_research} พบว่ายังมีช่องว่างในด้าน\ldots{} โครงงานนี้จึงมุ่งเน้นแก้ไขปัญหาดังกล่าว โดยพัฒนาระบบที่มีคุณสมบัติครบถ้วนทั้ง 4 ด้าน


\subsection{เทคโนโลยีและเครื่องมือที่เกี่ยวข้อง}

\subsubsection{ชื่อเครื่องมือ/เทคโนโลยี 1}

อธิบายคุณสมบัติและจุดเด่น\ldots{}

\subsubsection{ชื่อเครื่องมือ/เทคโนโลยี 2}

อธิบายคุณสมบัติและจุดเด่น\ldots{}
